First-order estimates predict the effect of removing a component of a network from the derivative of an output at the intact network. We tested them in the \NModels{} released FlyVis models, connectome-constrained recurrent networks of the fly visual system, where silencing each of the \NetTypes{} cell types can be simulated exactly. The outcome was the change in direction selectivity of the eight T4 and T5 motion-detector types, measured with moving edges at two speeds. In two studies, the analysis was fixed on four development models before the other \NConfModels{} were analysed. In the first, the estimate was a good screen and a poor estimate of effect size. At full silencing, a mean fraction of \PrecFOFull{} of the ten cell types it ranked highest changed direction selectivity by at least \LargeThreshold{}, against \PrecAnatFull{} for anatomical path weights through three synaptic steps (higher in \PrecWinsAnatFull{} of \NConfModels{} models). Its mean absolute error, \MaeFOFull{}, exceeded the \MaeZeroFull{} of predicting zero change, and \WrongSignFull\% of large changes had the wrong sign. The second study, with \PredPrimaryTotal{} predictions hashed before any run, asked why. It changed one thing at a time: each cell's output was replaced by its own average instead of zero, the index was built from mean instead of peak responses, or the second-order term of the Taylor expansion was added. It also removed one source of nonlinearity at a time from the original estimate. The nonlinearity that matters is in the readout. Applying the exact index to the linearised trajectories brings the error to \RatioCfReadFull{} times that of predicting zero change (95\% interval \RatioCICfReadFull{}), against \RatioRefFull{} for the original estimate, while applying the linearised index to the exactly simulated trajectories gives \RatioCfNetFull{}. Within the index, its ratio of a vector length to a sum of responses matters most. Mean replacement lowers the error ratio to \RatioAFull{}, the mean-response index changes nothing detectable (\RatioBFull{}), and the second-order term makes the estimates far worse (\RatioCFull{}). Exact silencing of the \SimulatedTypes{} cell types with a path to the readouts took \TimeExactFull{}~s per model, against \TimeGradient{}~s for the derivatives. In these models, first-order estimates are a useful screen for the cell types that matter; the effects of silencing them should be simulated, and an estimate of a ratio should apply the ratio exactly to linearised inputs.
